\documentclass[12pt,a4paper]{article}
\newcommand{\LabNumber}{5}
\newcommand{\LabTitle}{Criptografia cu chei publice:\\ RSA, ElGamal și Diffie–Hellman}
\newcommand{\LabTitleShort}{Lucrarea de laborator nr. 5 — Criptografie \& Securitate Informațională}
% ---------------------------------------------------------------------------
% Common preamble for the CS laboratory reports (compile with XeLaTeX)
% Strunga Daniel-Ioan, FAF-242 — https://github.com/don09dandidi/CS
% ---------------------------------------------------------------------------
\usepackage{fontspec}
\setmainfont{Latin Modern Roman}
\setsansfont{Latin Modern Sans}
\setmonofont{DejaVu Sans Mono}[Scale=0.85]
\usepackage[english]{babel}
\babelprovide[import, main]{romanian}
\usepackage[a4paper, margin=2.5cm, headheight=15pt]{geometry}
\usepackage{amsmath, amssymb}
\usepackage{booktabs, array, longtable, tabularx}
\usepackage{xcolor}
\usepackage{listings}
\usepackage{fvextra}
\usepackage[most]{tcolorbox}
\usepackage{fancyhdr}
\usepackage{tikz}
\usetikzlibrary{arrows.meta, positioning, shapes.geometric, fit, calc}
\usepackage{enumitem}
\usepackage{caption}
\usepackage{float}
\usepackage[hidelinks]{hyperref}
\usepackage{url}

\definecolor{codebg}{RGB}{248,248,250}
\definecolor{codeframe}{RGB}{200,200,210}
\definecolor{kw}{RGB}{0,0,160}
\definecolor{cmt}{RGB}{0,120,0}
\definecolor{str}{RGB}{160,30,30}
\definecolor{termbg}{RGB}{30,32,40}
\definecolor{termfg}{RGB}{225,228,235}
\definecolor{termbar}{RGB}{60,63,75}
\definecolor{boxblue}{RGB}{230,238,250}

\lstdefinestyle{py}{
  language=Python, basicstyle=\ttfamily\small, backgroundcolor=\color{codebg},
  frame=single, rulecolor=\color{codeframe}, numbers=left, numberstyle=\tiny\color{gray},
  keywordstyle=\color{kw}\bfseries, commentstyle=\color{cmt}\itshape, stringstyle=\color{str},
  breaklines=true, showstringspaces=false, columns=fullflexible, keepspaces=true, tabsize=4,
  extendedchars=true, inputencoding=utf8
}
\lstdefinestyle{sh}{
  language=bash, basicstyle=\ttfamily\small, backgroundcolor=\color{codebg},
  frame=single, rulecolor=\color{codeframe}, numbers=left, numberstyle=\tiny\color{gray},
  keywordstyle=\color{kw}\bfseries, commentstyle=\color{cmt}\itshape, stringstyle=\color{str},
  breaklines=true, showstringspaces=false, columns=fullflexible, keepspaces=true
}
\lstdefinestyle{wl}{
  basicstyle=\ttfamily\small, backgroundcolor=\color{codebg},
  frame=single, rulecolor=\color{codeframe}, numbers=left, numberstyle=\tiny\color{gray},
  morekeywords={RandomPrime,RandomInteger,PowerMod,Mod,GCD,FromDigits,IntegerDigits,
    ToCharacterCode,FromCharacterCode,IntegerLength,BitLength,Module,Print,Hash,Length},
  keywordstyle=\color{kw}\bfseries, commentstyle=\color{cmt}\itshape,
  morecomment=[s]{(*}{*)}, breaklines=true, columns=fullflexible, keepspaces=true
}

% Console session box (as in the Lab 2 report)
\newtcolorbox{terminal}[1][Sesiune Terminal Consolă]{
  enhanced, breakable, colback=termbg, colframe=termbar, coltext=termfg,
  fonttitle=\small\ttfamily\bfseries, title={\textbullet\textbullet\textbullet\quad #1},
  left=4pt, right=4pt, top=2pt, bottom=2pt, boxrule=0.6pt, arc=2pt
}
\DefineVerbatimEnvironment{termout}{Verbatim}{fontsize=\scriptsize, breaklines=true, breakanywhere=true}

\newtcolorbox{infobox}[1]{
  enhanced, breakable, colback=boxblue, colframe=kw!60, fonttitle=\bfseries,
  title={#1}, boxrule=0.6pt, arc=2pt
}

\newcommand{\bignum}[1]{\par\nobreak\vspace{1pt}{\noindent\ttfamily\footnotesize\seqsplit{#1}\par}\vspace{5pt}}
\usepackage{seqsplit}

\pagestyle{fancy}
\fancyhf{}
\renewcommand{\headrulewidth}{0.4pt}
\fancyhead[L]{\footnotesize\textit{\LabTitleShort}}
\fancyhead[R]{\footnotesize strunga.daniel-ioan}
\fancyfoot[C]{\thepage}

\newcommand{\makereporttitle}{%
\begin{titlepage}
  \centering
  {\large\bfseries UNIVERSITATEA TEHNICĂ A MOLDOVEI\par}
  \vspace{2pt}
  {\normalsize Facultatea Calculatoare, Informatică și Microelectronică\par}
  {\normalsize Departamentul Ingineria Software și Automatică\par}
  \vspace{5cm}
  {\LARGE\bfseries LUCRAREA DE LABORATOR NR. \LabNumber\par}
  \vspace{8pt}
  {\Large\bfseries \LabTitle\par}
  \vspace{10pt}
  {\normalsize Disciplina: Criptografie și Securitate Informațională\par}
  \vfill
  \begin{tabular*}{\textwidth}{@{}l@{\extracolsep{\fill}}r@{}}
    \textbf{Autor:} & \textbf{Verificat de:} \\
    Daniel-Ioan Strunga & Lector / Supervizor \\
    Student gr. FAF-242 & Maia Zaica \\
  \end{tabular*}
  \vspace{3cm}

  {\normalsize Chișinău — 2026\par}
\end{titlepage}
}


\begin{document}
\makereporttitle
\tableofcontents
\newpage

\section{Introduction and Laboratory Objectives}
Public-key (asymmetric) cryptography, introduced by Diffie and Hellman in 1976, removes the need
for a secret channel to distribute keys: every participant owns a key pair — a public key that
can be published and a private key that is never disclosed. Its security rests on
computationally hard number-theoretic problems: \textbf{integer factorisation} (RSA) and the
\textbf{discrete logarithm problem} in $\mathbb{Z}_p^*$ (ElGamal, Diffie–Hellman).

Objectives of the laboratory work:
\begin{enumerate}
  \item \textbf{Task 2.1} — generate RSA keys with a modulus of at least 2048 bits, encrypt and
        decrypt the message $m=$ \textit{„Nume Prenume”} $=$ \texttt{"Strunga Daniel"}.
  \item \textbf{Task 2.2} — generate ElGamal keys for the given 2048-bit prime $p$ and generator
        $g=2$, encrypt and decrypt the same message.
  \item \textbf{Task 3} — perform the Diffie–Hellman key exchange between Alice and Bob, who will
        use AES with a 256-bit key; the secrets $a$ and $b$ are chosen randomly according to the
        requirements of the algorithm.
  \item Comment in detail every step of the algorithms.
\end{enumerate}

\begin{infobox}{Tooling note}
The laboratory guide suggests Wolfram Alpha / Mathematica. All computations in this report were
performed with the script \texttt{lab5/public\_key\_crypto.py}, which uses only Python's
arbitrary-precision integers and the exact equivalents of the Wolfram functions
(\texttt{RandomPrime}, \texttt{RandomInteger}, \texttt{PowerMod}, \texttt{Mod}). The same steps
written in Wolfram Language are given in \texttt{lab5/wolfram/lab5.wl} and in
Appendix~\ref{app:wl}; random values will of course differ from run to run.
\end{infobox}

\section{Message Encoding}
As required, the message is converted to a decimal number through the hexadecimal representation
of its ASCII codes:

\begin{table}[H]
\centering
\caption{ASCII encoding of the message}
\small
\setlength{\tabcolsep}{3pt}
\begin{tabular}{l*{14}{c}}
\toprule
Char & S & t & r & u & n & g & a & \textvisiblespace & D & a & n & i & e & l \\
Dec & 83 & 116 & 114 & 117 & 110 & 103 & 97 & 32 & 68 & 97 & 110 & 105 & 101 & 108 \\
Hex & 53 & 74 & 72 & 75 & 6E & 67 & 61 & 20 & 44 & 61 & 6E & 69 & 65 & 6C \\
\bottomrule
\end{tabular}
\end{table}

\begin{align*}
  \text{hex}(m) &= \texttt{537472756E67612044616E69656C}\\
  m &= 1692665887211551163022180266632556 \quad (111\text{ bits})
\end{align*}
Since $m < n$ (RSA) and $m < p$ (ElGamal), the message is encrypted as a single block. Decoding is
the reverse process: the decimal number is written in base 256 and every digit is an ASCII code.

\section{Task 2.1 — RSA}
\subsection{Algorithm}
\begin{enumerate}
  \item Choose two large random primes $p$ and $q$ of 1024 bits each (the two most significant
        bits set, so that $n$ has exactly 2048 bits).
  \item Compute the modulus $n = p\cdot q$ and Euler's function $\varphi(n)=(p-1)(q-1)$.
  \item Choose the public exponent $e$ with $1<e<\varphi(n)$, $\gcd(e,\varphi(n))=1$; the
        standard value $e = 2^{16}+1 = 65537$ is used.
  \item Compute the private exponent $d = e^{-1} \bmod \varphi(n)$ (extended Euclid).
  \item Public key: $(e,n)$; private key: $(d,n)$.
  \item Encryption: $c = m^{e} \bmod n$; decryption: $m = c^{d} \bmod n$.
\end{enumerate}
Correctness follows from Euler's theorem: $c^d = m^{ed} = m^{1+t\varphi(n)} \equiv m \pmod n$.

\subsection{Key Generation}
\noindent\textbf{Prime} $p$ (1024 bits):
\bignum{176808871112039090984526464850735286439660389023386462388082435210186414461437480669924926133177454708045016279326372545269157284659126146165917432686200411796724137722531618215827249603125828229043112797177886064797073780214064020045599155225888975333919049719687450644337017021473196093008977146633192691913}
\noindent\textbf{Prime} $q$ (1024 bits):
\bignum{147544777097784200277020067528879447485998721355942058857431226508760019240145267226164837616924554294319508252767460643170882807222501248200899343651808553036405961156958405034586880688880740988104928633540855457200837417795714906394907658055887844083192750664298561387068898179994570602639067468375676501001}
\noindent\textbf{Modulus} $n=p\cdot q$ (\textbf{2048 bits}):
\bignum{26087225477136663758168126002811220423539097473803278654469747681066747888611776926290228230831132386896535866367977967409004018450730072346904226930218825458347163115556016306040293990379744792467377686123113862116316452575074350109224970044989314798246246101269374769666733480588650186197175941993049273323510473675817296320105588503821310555511765681007117646898532373588252650153552938525208351543951161603528563563429518347687047144838859831814031537663898619818690375808438782053179806172498229672515828160401029473283022686384908314078015442684648520906920619691628839966699172103608412132075724387280029104913}
\noindent\textbf{Euler's function} $\varphi(n)=(p-1)(q-1)$:
\bignum{26087225477136663758168126002811220423539097473803278654469747681066747888611776926290228230831132386896535866367977967409004018450730072346904226930218825458347163115556016306040293990379744792467377686123113862116316452575074350109224970044989314798246246101269374769666733480588650186197175941993049273323186120027607473028844041971441695821586106570627789125653018711869306216451970190629118587793849152601164039031335685159247007052957232437447214761325889654985560276928948758802765675880491660455367786729682287951285111488375129387637508629402871701489808819307642827935293256902140645436427679772271159912000}
\noindent\textbf{Public exponent:} $e = 65537$, \quad $\gcd(e,\varphi(n)) = 1$.\\[4pt]
\noindent\textbf{Private exponent} $d = e^{-1} \bmod \varphi(n)$:
\bignum{23429024892179789922364736384019189509881253920847569744166423222910839891600173613447618039252477549481105399370627510928583357828272904287142162334036496141917629336382677933949153363439910260461073641264024273468487573158922458658145360817218152794756790809384227681581614445793480954413225792602177200337662874997405255886060263155119806740661559830944062745423364334069540787557357118120135052855000492438651665064702184915210027742038996964092125267206013995190428190789340296235286706999860516086218358471128824732970541481823576615148025091444888009490648905131888081685817565772053897031336448169981967457473}
Check: $e\cdot d \bmod \varphi(n) = 1$.

\subsection{Encryption and Decryption}
\noindent\textbf{Ciphertext} $c = m^{e} \bmod n$:
\bignum{361751358558410855782783846576655976575107712956983301733248306032428806750642657575824892321788961099744934777720388441072464463434106551892893710632790682877308019218084378121784524001904905560545217704777565731586765792837559156464747914773724082256317493697625261844600358913086623798381104843572266654293601660392093449592967595524386370383188880734411872793963327629503657153510663053757703944292906904141686715171151422604427911916996408232961003143216379793102369786124543692417119888373432180137634250815380861682010885546335438258339452917580131726911678089049524147593311242182794275382485118054457161757}
\noindent\textbf{Decryption} $c^{d} \bmod n$:
\bignum{1692665887211551163022180266632556}
which is equal to $m$; converting it back to ASCII gives \texttt{"Strunga Daniel"}.

\section{Task 2.2 — ElGamal}
\subsection{Algorithm}
The public parameters are the given prime $p$ (2048 bits) and the generator $g=2$. The prime is a
\emph{safe prime} ($p=2q'+1$ with $q'$ prime — verified with the Miller–Rabin test), so the
discrete logarithm problem in $\mathbb{Z}_p^*$ has no small subgroups to exploit.
\begin{enumerate}
  \item \textbf{Key generation (receiver):} choose the private key $x$ randomly,
        $1 < x < p-1$, and publish $y = g^{x} \bmod p$. Public key: $(p,g,y)$.
  \item \textbf{Encryption (sender):} choose a fresh random ephemeral key $k$, $1<k<p-1$, and
        compute
        $c_1 = g^{k} \bmod p$, \; $c_2 = m\cdot y^{k} \bmod p$. The ciphertext is $(c_1,c_2)$.
  \item \textbf{Decryption (receiver):} the shared value is $s=c_1^{x}=g^{kx}$, and
        $m = c_2 \cdot s^{-1} \bmod p = c_2 \cdot c_1^{\,p-1-x} \bmod p$ (Fermat: $s^{-1}=s^{p-2}$,
        and $c_1^{p-1-x}=c_1^{-x}$).
\end{enumerate}
The ephemeral key $k$ must never be reused: from two ciphertexts with the same $k$ one obtains
$m_1/m_2 = c_2/c_2'$.

\subsection{Values}
\noindent\textbf{Prime} $p$ (given, 2048 bits) and $g = 2$:
\bignum{32317006071311007300153513477825163362488057133489075174588434139269806834136210002792056362640164685458556357935330816928829023080573472625273554742461245741026202527916572972862706300325263428213145766931414223654220941111348629991657478268034230553086349050635557712219187890332729569696129743856241741236237225197346402691855797767976823014625397933058015226858730761197532436467475855460715043896844940366130497697812854295958659597567051283852132784468522925504568272879113720098931873959143374175837826000278034973198552060607533234122603254684088120031105907484281003994966956119696956248629032338072839127039}
\noindent\textbf{Private key} $x$:
\bignum{25919002686015089196997653179970843150834060904206735136067689176923836682046556116865453758547975821853026443357016485951780236020323126156640097240177135448016435943616920695921518254802430583917018107831668171067767163959558718987457742498609667420362243617182097740572950280442786330235006799600528801421343646422551734743034037135226208403663913234983713635059288694711592237850852484085996870732664836062938879847236624612654301701702418710098331830327978456645938476105370784703538526021318421767784448515699519483737148593926776074671952197359720649113525693892564566311312211855001389191681857795763056793397}
\noindent\textbf{Public key} $y = g^{x} \bmod p$:
\bignum{28430938163656368085494262635577030208013298275637895066533828810679201616093866047507698485725500465157784853170964954364207086379297551081416080228346619857360226991450080799067528787822062186815894923129584416789993545096060739895462790607526271616062418927268533096829127995590800674324926471429219005910239910339374006715225863517212149342888203887060918700571499096240878140599462549725695931807359454288773520527401999196106165153891670473077873918871153744401335317836321299358019164815916787244809974457348113624632835642081250664930758829386044272880225742133533550225355660445594554201955008317732529743475}
\noindent\textbf{Ephemeral key} $k$:
\bignum{4530944111442528150505681461394463708644386694715735952301406439556345937514140637858436690327693653303450278277887222533017719853519071328025626208621379146805241729940827995393537858766917795047122183246230901246659505728253984389477738311028599855501528465925215427895928848103675651382737295636727705576046455888771902302798852799871790809324457954001909942435646376066819997687993886561357983673707846427317911388333296253146392387117216248133428250773212709978739957721580181937189382307176089604130693435568874891462557934598160450285752857776589923444837195439526971853026546646473375845189211714198202801644}
\noindent\textbf{Ciphertext} $c_1 = g^{k} \bmod p$:
\bignum{26284631159594511252032807001016336144146639419417864935737932610738838873919585514368283179636638858379965317689508610173682717273561399232349051899486333481810726150200708207151476350878582738473799896970438197439312562274274510036197727516890680034790798965533052117503125516129765881425869378993753183961744646298478078819922065718539792331105749470857855351831372028660826981894408790811729421877277267259003511249648297487295465235413027784989404028533191980481284912575706045923231534107537732546368352769785856417395401719881490023730956884094170275066090752083550215068083987380175893128492587176434335626807}
\noindent\textbf{Ciphertext} $c_2 = m \cdot y^{k} \bmod p$:
\bignum{11001169252029023168469692408091309376670827664202711081817940005797224622822809520079914636377991469495946970955815444130452657819424876412869323621509185940582195712651481848837946348463404613192948715907658656974622330342819733657464084697229964907165107692327211544735179269245129354296593843031406225236846954369481674010366952548914283506749101048679752294451089810875131831540113450164234152891486347617371668303918092271624614903876917373784569749399019631907867397948581758379812529878172699867192747007354258259483643634411971550520284495661970648071744721216297770780615923740356270604155851727373726132331}
\noindent\textbf{Decryption} $m = c_2 \cdot c_1^{\,p-1-x} \bmod p$:
\bignum{1692665887211551163022180266632556}
The decrypted number equals $m$ and decodes to \texttt{"Strunga Daniel"}.

\section{Task 3 — Diffie–Hellman Key Exchange for AES-256}
\subsection{Protocol}
Alice and Bob agree publicly on $p$ (2048 bits) and $g=2$.
\begin{enumerate}
  \item Alice chooses a random secret $a$, $1<a<p-1$, and sends $A = g^{a} \bmod p$ to Bob.
  \item Bob chooses a random secret $b$, $1<b<p-1$, and sends $B = g^{b} \bmod p$ to Alice.
  \item Alice computes $K_A = B^{a} \bmod p$, Bob computes $K_B = A^{b} \bmod p$; both equal
        $g^{ab} \bmod p$. An eavesdropper sees only $p,g,A,B$ and would need to solve the
        discrete logarithm (or the Diffie–Hellman problem) to obtain $K$.
  \item The shared secret has 2048 bits while AES-256 needs exactly 256 bits, so a key-derivation
        step is applied: $K_{AES} = \mathrm{SHA\text{-}256}(K)$, where $K$ is written as 256
        big-endian bytes. Both parties obtain the same 256-bit key.
\end{enumerate}
The exchange as such is not authenticated; in practice the values $A$ and $B$ are signed (or
certified, see the PKI laboratory) to prevent man-in-the-middle attacks.

\subsection{Values}
\noindent\textbf{Alice's secret} $a$:
\bignum{23792718007091579695681404510124568152371853167751596220032554300185142679410360006580727950943368932083278253342028392355795936792658083873352519425086939222988198006741053254685468457992373327386577281749450898065290845060210853390001083775573921938784911728145075374311689047963089438256798954678321586190638917047679252624971166252102144213574409997269828587193934757957761567813854059327839677220649744888415704456238795116151843795111736748847453774008829230870664481233191569578227462624168093823318672479486178419541711457094189485487557546428882366295829779224060848078470109345007231984000348063048668701346}
\noindent\textbf{Bob's secret} $b$:
\bignum{30013154077667882323168009006877076022355367825449596763356465993925503532833149117692248724787274159469923116351160886674927629163828831281214455015419391548571095752769858375436629368683720457626370058422906968344370232307740406914148065989576378150867776347524624869413920903627393428600688295520459456021062714551144681390686522271751247877651087388324225272771814779296310693586180118213506292065705817046162072275951257368051316784629116134448758172488806975658023056301131947855258377486420288919633673917749613041159904822183515596033597411711037435356315000468998049804607191056217949777683208954325040943151}
\noindent\textbf{Alice sends} $A = g^{a} \bmod p$:
\bignum{7311244260793985346449738252913125163428217206109081186839667044362971263738288453510227084305807329816778591187318008933484363670240067973617720961235769476559114928293192920849991643360153037067643788895565644907803414984542181355631526071817434846553346231777306192584622181266610239574750089290961935104362791116238137227250119929423132214364525694497975488442395509374085446708001087565502358072930664205535672234754173120826871833646734629049659178626804245294961086086710333194920756046808103432033497077008796713478358541855503483883714845425622427396667372423027106318323791945150067548721921379310688461158}
\noindent\textbf{Bob sends} $B = g^{b} \bmod p$:
\bignum{5261837725609434206087303441633177365634439994251374366560684401033329057675780720317728148405147452496232611281954057155317958286065559332016596222299705102291156764093477069877830200260565533254366999395109429074823286544204424632267533648349234735902079515612480196409007754879845618468613410675740937142221772571212371838821611514759335680783974523009981584510865841504109162621633243877518144291936567968223147781503543120651610621702247896689556995785118694846231237583442377106854618641703175991079208205033226386567068856368119018054524209346162246300424586017883481799806217310907675427809060834063211360391}
\noindent\textbf{Alice computes} $K_A = B^{a} \bmod p$:
\bignum{15454418984412565387331444407078855129471671681186594013469844569473635427341839006923362821161887917499157945701933519209178515128835213352828537713105604583699401247028906685082167886957883063929537020604765605476851059031006218803072359488538374185045873630618285469236579209052016622458087042767131386546655425004681012371038139766536259696445317814051346902737270509549654731370009427188371645870875273392769891706831080258403069157511794356534684819973472927791932539169593209563024958735745372114729896902155459762203787629868832744934614212637656978563397092915604746387975796186110503858081805236624537262094}
\noindent\textbf{Bob computes} $K_B = A^{b} \bmod p$:
\bignum{15454418984412565387331444407078855129471671681186594013469844569473635427341839006923362821161887917499157945701933519209178515128835213352828537713105604583699401247028906685082167886957883063929537020604765605476851059031006218803072359488538374185045873630618285469236579209052016622458087042767131386546655425004681012371038139766536259696445317814051346902737270509549654731370009427188371645870875273392769891706831080258403069157511794356534684819973472927791932539169593209563024958735745372114729896902155459762203787629868832744934614212637656978563397092915604746387975796186110503858081805236624537262094}
$K_A = K_B$: \textbf{True}.\\[4pt]
\noindent\textbf{AES-256 key} $K_{AES}=\mathrm{SHA\text{-}256}(K)$ (256 bits, hex):
\bignum{1CCAE0F19718B4C8F51D665B304501586B11673C9FEE7860A71F62ACB2678A05}

\subsection{Using the Key}
To demonstrate that the key is usable, Alice encrypts the message with AES-256-GCM and Bob decrypts it:
\begin{terminal}
\begin{termout}
nonce      = 1F0722E548E79D9F1CA12F7A
ciphertext = F960501468D3B3811564FDEF397048A4186B88A55EE254E0E5B05C444C56
decrypted  = Strunga Daniel
\end{termout}
\end{terminal}

\section{Source Code and Repository}
\begin{center}
\url{https://github.com/don09dandidi/CS} \quad (directory \texttt{lab5/})
\end{center}
\begin{lstlisting}[style=py, caption={Core of the three tasks (excerpt from \texttt{public\_key\_crypto.py})}]
# RSA
n, phi = p * q, (p - 1) * (q - 1)
d = pow(e, -1, phi)                    # PowerMod[e, -1, phi]
c = pow(m, e, n); m2 = pow(c, d, n)    # PowerMod[m, e, n], PowerMod[c, d, n]

# ElGamal
x = random_in(2, P - 2); y = pow(G, x, P)
k = random_in(2, P - 2)
c1, c2 = pow(G, k, P), (m * pow(y, k, P)) % P
m2 = (c2 * pow(c1, P - 1 - x, P)) % P

# Diffie-Hellman
a, b = random_in(2, P - 2), random_in(2, P - 2)
A, B = pow(G, a, P), pow(G, b, P)
k_alice, k_bob = pow(B, a, P), pow(A, b, P)
aes_key = hashlib.sha256(k_alice.to_bytes(256, "big")).digest()
\end{lstlisting}

\subsection{Build and Execution Instructions}
\begin{lstlisting}[style=sh]
git clone https://github.com/don09dandidi/CS.git
cd CS/lab5
pip install cryptography          # optional, only for the AES-256-GCM demo
python3 public_key_crypto.py --save   # prints all values, saves results.json
\end{lstlisting}

\section{Conclusion}
\begin{enumerate}
  \item \textbf{RSA:} keys with a 2048-bit modulus were generated; the message
        \texttt{"Strunga Daniel"} was encrypted with $(e,n)$ and correctly recovered with $(d,n)$.
        Security relies on the difficulty of factoring $n$ — knowing $p,q$ immediately yields
        $\varphi(n)$ and $d$.
  \item \textbf{ElGamal:} with the 2048-bit safe prime and $g=2$ the message was encrypted into the
        pair $(c_1,c_2)$ and decrypted correctly. ElGamal is probabilistic (a fresh $k$ gives a
        different ciphertext each time) and the ciphertext is twice as long as the message.
  \item \textbf{Diffie–Hellman:} Alice and Bob obtained the same 2048-bit secret without
        transmitting it, and derived from it a 256-bit AES key with SHA-256, which was then used
        for authenticated encryption with AES-256-GCM.
  \item In practice textbook RSA/ElGamal are combined with padding schemes (OAEP) and authenticated
        key exchange; they are used to transport symmetric keys rather than to encrypt bulk data.
\end{enumerate}

\appendix
\section{Wolfram Language Code}\label{app:wl}
\lstinputlisting[style=wl]{wolfram/lab5.wl}

\begin{thebibliography}{9}
\bibitem{rsa} R.\ Rivest, A.\ Shamir, L.\ Adleman, \emph{A Method for Obtaining Digital Signatures
  and Public-Key Cryptosystems}, Communications of the ACM, 21(2), 1978.
\bibitem{elg} T.\ ElGamal, \emph{A Public Key Cryptosystem and a Signature Scheme Based on Discrete
  Logarithms}, IEEE Trans.\ Inf.\ Theory, 31(4), 1985.
\bibitem{dh} W.\ Diffie, M.\ Hellman, \emph{New Directions in Cryptography}, IEEE Trans.\ Inf.\
  Theory, 22(6), 1976.
\bibitem{rfc7919} D.\ Gillmor, \emph{RFC 7919: Negotiated Finite Field Diffie-Hellman Ephemeral
  Parameters for TLS}, 2016.
\bibitem{stallings} William Stallings, \emph{Cryptography and Network Security: Principles and
  Practice}, 7th Edition, Pearson, 2017.
\end{thebibliography}
\end{document}
